\documentclass[11pt]{article}
\usepackage{fontspec,amsmath,amssymb,amsthm,hyperref}
\usepackage[margin=1in]{geometry}
\setlength{\emergencystretch}{3em}
\title{Commute the actual gradient and Laplacian under C\ensuremath{{}^{3}} regularity}
\author{LangevinLean mathematical exposition}
\date{}
\begin{document}\maketitle
\noindent Source module: \texttt{AverageLMC/LaplacianGradient.lean}. Compile this proof document with XeLaTeX or LuaLaTeX. The accompanying HTML pairs the mathematical argument with the exact Lean statements and records the scope of the formalization.


The third-order analysis needs the identity Δ∇f=∇Δf for the true derivatives. The proof uses Schwarz symmetry twice, and requires three continuous derivatives rather than analyticity or a Lipschitz Hessian.



Three continuous derivatives suffice for the directional Schwarz identities and the vector Laplacian-gradient commutation used by the diffusion generator:


\[D_uD_uD_vf=D_vD_uD_uf\qquad\bigl(f\in C^3(\mathbb R^p)\bigr).\]
\[\Delta(\nabla f)=\nabla(\Delta f).\]
\section{Fixed directional derivatives and Schwarz symmetry}


A directional derivative is the actual Fréchet derivative applied to a fixed vector. A continuously differentiable derivative map loses one order of regularity, and evaluation at a fixed vector preserves the remaining regularity. Differentiating that evaluation identifies a second directional derivative with the actual second Fréchet derivative applied to the two directions. Schwarz symmetry for C\ensuremath{{}^{2}} functions then commutes the two directional derivatives as functions of the base point. These statements are proved for arbitrary real normed spaces, before choosing Euclidean coordinates.



\noindent Lean declarations: \texttt{directionalDerivative}, \texttt{contDiff\_directionalDerivative}, \texttt{directionalDerivative\_two}, \texttt{directionalDerivative\_commute}.

\section{Commute the needed three derivatives without analyticity}


For a C\ensuremath{{}^{3}} potential, first commute the two inner directional derivatives using C\ensuremath{{}^{2}} symmetry of f. Then commute the outer pair using C\ensuremath{{}^{2}} symmetry of one first directional derivative. This gives DᵤDᵤDᵥf=DᵥDᵤDᵤf. The argument establishes exactly the third-order permutation needed for the Laplacian and avoids any appeal to an analytic-only iterated-derivative permutation theorem.



\noindent Lean declarations: \texttt{directionalDerivative\_three}.

\section{Expand the canonical trace and differentiate the finite sum}


The canonical orthonormal-basis formula for the Laplacian expresses it as the sum of second directional derivatives along the coordinate vectors. The jth coordinate of the actual gradient equals the directional derivative along e\_j by the gradient–Fréchet-derivative identity. A finite sum of differentiable scalar functions can be differentiated term by term. Finally, the gradient is C\ensuremath{{}^{2}} because it is the continuous Riesz inverse applied to the C\ensuremath{{}^{2}} derivative map of a C\ensuremath{{}^{3}} potential.



\noindent Lean declarations: \texttt{laplacian\_eq\_directionalDerivative}, \texttt{gradient\_coordinate\_directionalDerivative}, \texttt{directionalDerivative\_sum}, \texttt{contDiff\_gradient\_of\_three}.

\section{Identify the vector Laplacian coordinate by coordinate}


Continuous linear projection onto coordinate j commutes with the vector Laplacian. Its scalar observable is Dⱼf, whose Laplacian is the sum of DᵢDᵢDⱼf. The jth coordinate of ∇Δf is Dⱼ of the finite Laplacian sum, namely the sum of DⱼDᵢDᵢf. Apply the proved third-derivative commutation term by term, then use equality of all Euclidean coordinates to obtain the full vector identity.



\noindent Lean declarations: \texttt{laplacian\_gradient\_eq\_gradient\_laplacian}.

\end{document}
